% SHARD-ID: AQM-70-QSA-AND-OR-MANY-PARTICLE-GRAMMAR
% SHARD-TITLE: AND/OR Many-Particle Grammar
% SHARD-SUMMARY: Defines a finite carrier grammar whose AND operation is tensor/coexistence and whose OR operation is direct-sum/alternative.
% SHARD-SUMMARY: Identifies finite particle-count sectors by a recursive grading on finite tensor/direct-sum expressions.
% SHARD-SUMMARY: Separates this elementary grammar from symmetric, antisymmetric, para, Fock, and fusion-category enhancements.
% SHARD-KEYWORDS: quantum-system association, finite grammar, tensor product, direct sum, particle sectors, Fock caveat

\section{AND/OR Many-Particle Grammar}
\label{sec:qsa-and-or-many-particle-grammar}

\subsection{Status and Local Anchors}

This is QSA Step 06 from
\path{docs/quantum_system_associations/PLAN.md}.  It is a local
finite-dimensional construction using \texttt{CONVENTIONS.md} conventions
\texttt{(ap)}, \texttt{(ar)}, \texttt{(as)}, and \texttt{(at)}, with
AQM-64--AQM-66 as the planning and catalogue anchors.  AQM-68 supplies the
finite tensor/coexistence operation, and AQM-69 supplies the finite
direct-sum/alternative operation.

The shard uses no external source and proves no canonical bosonic, fermionic,
para, Fock, or fusion-category theorem.  AQM-23 is used only as the local
caution that direct sums and tensor products live in different layers unless a
workflow explicitly relates them.  AQM-19 is used only as a boundary marker
for later symmetric and antisymmetric field enhancements; no canonical-field
equivalence from that shard is imported here.

\subsection{Finite Grammar}

The input is a finite list of already chosen finite-dimensional complex
carriers.  These may be the site Hilbert spaces \(\mathcal H_x\) from AQM-68
or the internal carriers \(K_x\) from AQM-69.  The bare finite set still
supplies only labels; the carriers are extra data.

A grammar expression is built by the rules
\[
  E ::= 0 \mid \mathbf 1 \mid K \mid
        \operatorname{AND}(E,F) \mid \operatorname{OR}(E,F),
\]
where \(K\) ranges over the chosen carrier atoms, \(0\) is the zero carrier,
and \(\mathbf 1\) is the tensor unit \(\mathbb C\).  Its carrier
\(\mathcal V(E)\) is defined recursively by
\[
\begin{aligned}
  \mathcal V(0)&=0,\\
  \mathcal V(\mathbf 1)&=\mathbb C,\\
  \mathcal V(K)&=K,\\
  \mathcal V(\operatorname{AND}(E,F))
    &=\mathcal V(E)\otimes_{\mathbb C}\mathcal V(F),\\
  \mathcal V(\operatorname{OR}(E,F))
    &=\mathcal V(E)\oplus\mathcal V(F).
\end{aligned}
\]

Thus \(\operatorname{AND}\) means independent coexistence in the finite tensor
product sense of AQM-68, while \(\operatorname{OR}\) means alternatives in the
finite direct-sum sense of AQM-69.  Parentheses are part of the written
grammar.  Associativity and distributivity may be used as finite vector-space
isomorphisms after the expression is fixed, but they are not new structure
coming from the bare set.

If the atoms are Hilbert spaces, the same recursion uses the finite Hilbert
tensor product and finite Hilbert direct sum.  If the atoms are only complex
carriers, the output is only a complex vector space.  No observable algebra is
selected by this grammar alone; \(\operatorname{End}(\mathcal V(E))\), a
block-preserving direct-sum algebra, and a tensor-site local algebra would be
different later choices.

\subsection{Particle-Count Sectors}

When every carrier atom is declared to be one particle and \(\mathbf 1\) is
declared to be the zero-particle carrier, the finite grammar has a recursive
particle-count grading.  Write \(\mathcal V_n(E)\) for the degree-\(n\)
sector.  The definitions are
\[
\begin{array}{rclcrcl}
  \mathcal V_0(\mathbf 1)&=&\mathbb C, &&
  \mathcal V_n(\mathbf 1)&=&0 \quad(n\neq0),\\
  \mathcal V_1(K)&=&K, &&
  \mathcal V_n(K)&=&0 \quad(n\neq1),\\
  \mathcal V_n(0)&=&0 &&
  \text{for all }n.
\end{array}
\]
The operations are
\[
  \mathcal V_n(\operatorname{OR}(E,F))
  =
  \mathcal V_n(E)\oplus\mathcal V_n(F)
\]
and
\[
  \mathcal V_n(\operatorname{AND}(E,F))
  =
  \bigoplus_{i+j=n}\mathcal V_i(E)\otimes_{\mathbb C}\mathcal V_j(F).
\]
Because the parse tree is finite, only finitely many sectors are nonzero, and
\[
  \mathcal V(E)=\bigoplus_{n\geq0}\mathcal V_n(E)
\]
is a finite direct sum.  The integer \(n\) is a grammar count of declared
one-particle atoms.  It is not a Hamiltonian, a dynamics, a statistics rule,
or an infinite-particle completion.

\subsection{Basic Finite Sectors}

The AQM-69 one-particle alternatives carrier is recovered by the expression
\[
  P_X=\operatorname{OR}_{x\in X} K_x,
  \qquad
  \mathcal V(P_X)=\bigoplus_{x\in X}K_x .
\]
If all \(K_x\) are declared one-particle atoms, then
\(\mathcal V_1(P_X)=\mathcal V(P_X)\) and all other sectors vanish.

A finite site-option expression is
\[
  G_X=\operatorname{AND}_{x\in X}(\mathbf 1\operatorname{OR}K_x).
\]
After choosing a display order for the finite tensor product, finite
distributivity gives
\[
  \mathcal V_n(G_X)
  \cong
  \bigoplus_{\substack{S\subset X\\ |S|=n}}
  \bigotimes_{x\in S}K_x .
\]
This expression allows each displayed site factor to choose either the
zero-particle alternative \(\mathbf 1\) or the one-particle alternative
\(K_x\).  That site-option rule is part of the expression, not a consequence
of the bare finite set.

By contrast, if
\[
  P_X^{\operatorname{AND} m}
  =
  \underbrace{P_X\operatorname{AND}\cdots\operatorname{AND}P_X}_{m
  \text{ factors}},
\]
then
\[
  \mathcal V(P_X^{\operatorname{AND} m})
  \cong
  \bigoplus_{(x_1,\ldots,x_m)\in X^m}
  K_{x_1}\otimes\cdots\otimes K_{x_m},
\]
and the whole carrier lies in degree \(m\).  The \(m\) tensor factors are
grammar slots.  They are not symmetrized, antisymmetrized, or quotiented, and
the expression does not forbid two slots from choosing the same site.

A finite truncated many-sector expression can be written as
\[
  T_{\leq M}(X)=
  \mathbf 1
  \operatorname{OR} P_X
  \operatorname{OR} P_X^{\operatorname{AND}2}
  \operatorname{OR}\cdots
  \operatorname{OR} P_X^{\operatorname{AND}M}.
\]
It has sectors \(0,\ldots,M\) by construction.  This is still a finite
direct-sum/tensor expression, not an infinite Fock completion.

\subsection{Boundary With Enhancements}

The grammar above is elementary vector-space bookkeeping.  It does not choose
statistics.

For symmetric or antisymmetric sectors, one would need extra data and a new
rule: an action of the permutation group on tensor slots, plus a choice of
invariants, coinvariants, projectors, or quotients and compatible Hilbert
structure.  None of that is supplied by \(\operatorname{AND}\) and
\(\operatorname{OR}\) alone.  AQM-19 belongs to that later canonical-field
boundary; this shard does not claim that \(T_{\leq M}(X)\) is a bosonic or
fermionic Fock construction.

Para-statistics would likewise require a separately recorded permutation-sector
or algebraic selection rule.  Fusion-category and Levin-Wen/string-net
variants require their own objects, tensor product, direct-sum or fusion
rules, associativity data, and any pivotal, spherical, or braiding conventions
needed by the model.  Those choices are not hidden in this finite carrier
grammar and are deferred to later QSA shards.

Infinite-particle constructions are also out of scope.  The expression
\(\bigoplus_{m\geq0}P_X^{\operatorname{AND}m}\), any Hilbert direct-sum
completion of it, and any infinite tensor product require a separate
completion convention.  Step 06 only records finite expressions and their
finite particle-count sectors.

\subsection{Conclusion}

QSA Step 06 adds a syntax for combining the AQM-68 tensor-site operation and
the AQM-69 direct-sum particle operation:
\[
  \operatorname{AND}=\otimes_{\mathbb C},
  \qquad
  \operatorname{OR}=\oplus.
\]
With one-particle atoms and the vacuum unit \(\mathbf 1=\mathbb C\), this
syntax identifies finite particle-count sectors by recursion.  It remains
separate from statistics, Fock completions, and fusion-category enhancements.
